%Data institusi

%\input{Dozentdaten}

%\renewcommand{\fachbereich}{Fachbereich}

%\renewcommand{\dozent}{Prof. Dr. . }

%Data ujian

\renewcommand{\klausurgebiet}{ }

\renewcommand{\klausurtyp}{ }

\renewcommand{\klausurdatum}{ . 20}

\klausurvorspann {\fachbereich} {\klausurdatum} {\dozent} {\klausurgebiet} {\klausurtyp}

%Data untuk tabel nilai berikut

\renewcommand{\aeins}{  3  }

\renewcommand{\azwei}{  3   }

\renewcommand{\adrei}{  3  }

\renewcommand{\avier}{  0  }

\renewcommand{\afuenf}{  2  }

\renewcommand{\asechs}{  0  }

\renewcommand{\asieben}{  3  }

\renewcommand{\aacht}{  7  }

\renewcommand{\aneun}{  0  }

\renewcommand{\azehn}{  7  }

\renewcommand{\aelf}{  4  }

\renewcommand{\azwoelf}{  3  }

\renewcommand{\adreizehn}{  3   }

\renewcommand{\avierzehn}{  7  }

\renewcommand{\afuenfzehn}{  2  }

\renewcommand{\asechzehn}{  9  }

\renewcommand{\asiebzehn}{  7  }

\renewcommand{\aachtzehn}{  63  }

\renewcommand{\aneunzehn}{    }

\renewcommand{\azwanzig}{    }

\renewcommand{\aeinundzwanzig}{    }

\renewcommand{\azweiundzwanzig}{    }

\renewcommand{\adreiundzwanzig}{    }

\renewcommand{\avierundzwanzig}{    }

\renewcommand{\afuenfundzwanzig}{    }

\renewcommand{\asechsundzwanzig}{    }

\punktetabellesiebzehn

\klausurnote

\newpage

\setcounter{section}{0}



\inputaufgabeklausurloesung
{3}

{

Definisikan istilah-istilah berikut
\zusatzklammer {yang dicetak miring} {} {}.
\aufzaehlungsechs{\stichwort {Pemetaan Gauss} {}
dari suatu hiperpermukaan diferensiabel

\mavergleichskette
{\vergleichskette
{ Y
}
{ \subseteq }{ \R^n
}
{  }{
}
{    }{
}
{   }{
}
}
{}{}{.}

}{Suatu
\stichwort {peta topologis} {}
pada suatu
\definitionsverweis {manifold topologis}{}{}
$M$.

}{Suatu $C^1$-\stichwort {manifold diferensiabel} {} $M$.

}{\stichwort {Pangkat eksterior} {} ke-$n$ dari suatu
$K$-\definitionsverweis {ruang vektor}{}{}
$V$
\zusatzklammer {cukup tuliskan simbolnya} {} {.}

}{Suatu
\stichwort {isometri lokal} {}
\maabb {\varphi} {L} {M
} {}
antara manifold-manifold Riemann.

}{Suatu koneksi
\stichwort {terintegralkan secara lokal} {}
pada suatu bundel vektor diferensiabel
\maabb {p} {E} {X
} {}
di atas suatu manifold diferensiabel $X$.
}

}
{

\aufzaehlungsechs{Misalkan suatu
\definitionsverweis {orientasi}{}{}
pada $Y$ telah ditetapkan. Maka pemetaan
\maabbdisp {} {Y} { S^{n-1}
} {,}
yang memetakan setiap titik

\mavergleichskette
{\vergleichskette
{ P
}
{ \in }{ Y
}
{  }{
}
{    }{
}
{   }{
}
}
{}{}{}
ke vektor normal satuan yang ditentukan oleh orientasi tersebut disebut
pemetaan Gauss
dari $Y$.
}{Suatu peta topologis adalah setiap
\definitionsverweis {homeomorfisme}{}{}
\maabbdisp {\varphi} {U} {V
} {,}
dengan

\mavergleichskette
{\vergleichskette
{ U
}
{ \subseteq }{ M
}
{  }{
}
{    }{
}
{   }{
}
}
{}{}{}
dan

\mavergleichskette
{\vergleichskette
{ V
}
{ \subseteq }{ \R^n
}
{  }{
}
{    }{
}
{   }{
}
}
{}{}{}
\definitionsverweis {terbuka}{}{.}
}{Suatu
\definitionsverweis {ruang Hausdorff}{}{}
\definitionsverweis {topologis}{}{}
$M$ bersama suatu
\definitionsverweis {penutup terbuka}{}{}

\mathl{M= \bigcup_{i \in I} U_i}{} dan
\definitionsverweis {peta-peta}{}{}
\maabbdisp {\alpha_i} {U_i } { V_i
} {}
dengan
\mathl{V_i \subseteq \R^n}{} terbuka, sedemikian rupa sehingga
\definitionsverweis {pemetaan-pemetaan transisi}{}{}
\maabbdisp {\alpha_j \circ (\alpha_i)^{-1}} {V_i \cap \alpha_i(U_i \cap U_j)   } { V_j \cap \alpha_j(U_i \cap U_j)
} {}
$C^1$-\definitionsverweis {difeomorfisme}{}{}
untuk semua
\mathl{i,j \in  I}{}. Struktur ini disebut manifold diferensiabel.
}{Ruang vektor-$K$
\mathl{\bigwedge^n V}{} disebut pangkat eksterior ke-$n$ dari $V$.
}{Suatu
\definitionsverweis {pemetaan diferensiabel}{}{}
\maabb {\varphi} {L} {M
} {}
disebut
isometri lokal
jika untuk setiap titik

\mavergleichskette
{\vergleichskette
{ P
}
{ \in }{ L
}
{  }{
}
{    }{
}
{   }{
}
}
{}{}{}
pemetaan
\definitionsverweis {tangen}{}{}
\maabbdisp {T_P \varphi} {T_PL} { T_{\varphi(P) } M
} {}
merupakan suatu
\definitionsverweis {isometri}{}{}
terhadap hasil kali skalar yang diberikan.
}{Koneksi tersebut disebut
terintegralkan secara lokal
jika untuk setiap titik

\mavergleichskette
{\vergleichskette
{ e
}
{ \in }{ E
}
{  }{
}
{    }{
}
{   }{
}
}
{}{}{}
terdapat suatu

\mavergleichskette
{\vergleichskette
{ p(e)
}
{ \in }{ U
}
{ \subseteq }{ X
}
{    }{
}
{   }{
}
}
{}{}{}
\definitionsverweis {penampang horizontal}{}{}
\maabbdisp {s} {U} {E \vert_U
} {}
yang didefinisikan pada lingkungan terbuka tersebut dan melalui $e$.
}

 }



\inputaufgabeklausurloesung
{3}

{

Rumuskan teorema-teorema berikut.
\aufzaehlungdrei{\stichwort {Teorema isometri untuk transpor paralel pada suatu hiperpermukaan} {}

\mavergleichskette
{\vergleichskette
{  Y
}
{ \subseteq }{  \R^n
}
{  }{
}
{    }{
}
{   }{
}
}
{}{}{.}}{\stichwort {Theorema egregium} {.}}{Teorema tentang partisi kesatuan.}

}
{

\aufzaehlungdrei{\faktsituation {Misalkan

\mavergleichskette
{\vergleichskette
{ Y
}
{ \subseteq }{ W
}
{ \subseteq }{ \R^n
}
{    }{
}
{   }{
}
}
{}{}{,}
$W$
\definitionsverweis {terbuka}{}{,}
dan merupakan suatu
\definitionsverweis {hiperpermukaan diferensiabel}{}{.}
Misalkan
\maabb {\gamma} {[a,b]} {Y
} {}
suatu
\definitionsverweis {kurva diferensiabel}{}{}
dengan

\mavergleichskette
{\vergleichskette
{ \gamma(a)
}
{ = }{ P
}
{  }{
}
{    }{
}
{   }{
}
}
{}{}{}
dan

\mavergleichskette
{\vergleichskette
{ \gamma(b)
}
{ = }{ Q
}
{  }{
}
{    }{
}
{   }{
}
}
{}{}{.}}
\faktfolgerung {Maka
\definitionsverweis {transpor paralel}{}{}
sepanjang $\gamma$ merupakan suatu
\definitionsverweis {isometri}{}{}
\maabbdisp {} {T_PY} {T_QY
} {.}}
\faktzusatz {}
\faktzusatz {}}{\faktsituation {Misalkan

\mavergleichskette
{\vergleichskette
{ Y
}
{ \subseteq }{ \R^3
}
{  }{
}
{    }{
}
{   }{
}
}
{}{}{}
suatu
\definitionsverweis {permukaan berorientasi}{}{}
dan misalkan
\maabbdisp {\varphi} { V } { Y
} {,}

\mavergleichskette
{\vergleichskette
{ V
}
{ \subseteq }{ \R^2
}
{  }{
}
{    }{
}
{   }{
}
}
{}{}{,}
suatu parametrisasi lokal $Y$ yang dua kali terdiferensialkan, dengan parameter $u,v$. Misalkan
\mathl{\left(  g_{ij}  \right)}{}
\definitionsverweis {matriks fundamental pertama}{}{}
pada $V$.}
\faktfolgerung {Maka, dengan menggunakan
\definitionsverweis {simbol-simbol Christoffel}{}{,}
\definitionsverweis {kelengkungan Gauss}{}{}
memenuhi hubungan

\mavergleichskettedisphandlinks
{\vergleichskettedisphandlinks
{ K
}
{ =} { { \frac{   1  }{ g_{11}  } }  \left(  \partial_2 \Gamma_{11}^2 - \partial_1 \Gamma_{12}^2 + \Gamma^1_{11} \Gamma^2_{12}  + \Gamma_{11}^2 \Gamma_{22}^2 - \Gamma_{12}^1 \Gamma_{11}^2-  \left(  \Gamma_{12}^2  \right)^2  \right)
}
{ } {
}
{ } {
}
{ } {
}
}
{}{}{.}}
\faktzusatz {}
\faktzusatz {}}{Misalkan $M$ suatu
\definitionsverweis {manifold diferensiabel}{}{}
dengan
\definitionsverweis {basis topologi yang terhitung}{}{.}
Maka, untuk setiap penutup terbuka, terdapat suatu
\definitionsverweis {partisi kesatuan}{}{}
\definitionsverweis {yang terdiferensialkan secara kontinu}{}{}
dan tersubordinasi pada penutup tersebut.}

 }



\inputaufgabeklausurloesung
{3}

{

Misalkan
\maabbdisp {f,g} {\R} { \R^n
} {}
\definitionsverweis {kurva-kurva diferensiabel}{}{.}
Hitung
\definitionsverweis {turunan}{}{}
dari fungsi
\maabbeledisp {} {\R} {\R
} { t } { \left\langle f(t) , g(t)  \right\rangle
} {.}
Rumuskan hasilnya dengan menggunakan hasil kali dalam.

}
{

Misalkan $f_i(t)$ dan $g_i(t)$ masing-masing fungsi komponen dari
\mathkor {} {f} {dan} {g} {.}
Fungsi yang dimaksud adalah

\mavergleichskettedisp
{\vergleichskette
{  \left\langle f(t) , g(t)  \right\rangle
}
{ =} { \sum_{i = 1}^n  f_i(t) g_i(t)
}
{ } {
}
{ } {
}
{ } {
}
}
{}{}{}
dengan turunan
\zusatzklammer {aturan hasil kali diterapkan pada setiap suku} {} {}

\mathdisp {\sum_{i = 1}^n  f'_i(t) g_i(t)  + \sum_{i = 1}^n  f_i(t) g'_i(t)} { . }

Dengan hasil kali skalar, bentuk ini dapat ditulis sebagai

\mathdisp {\left\langle f'(t) ,  g(t)  \right\rangle + \left\langle f(t) ,  g'(t)  \right\rangle} {  }

.

 }



\inputaufgabeklausurloesung
{0}

{

}
{  }



\inputaufgabeklausurloesung
{2}

{

Misalkan

\mavergleichskette
{\vergleichskette
{ Y
}
{ \subseteq }{ W
}
{ \subseteq }{ \R^n
}
{    }{
}
{   }{
}
}
{}{}{,}
$W$
\definitionsverweis {terbuka}{}{,}
suatu
\definitionsverweis {hiperpermukaan diferensiabel}{}{},
dan
\maabb {\gamma} {I =[a,b]} {Y
} {}
suatu
\definitionsverweis {kurva geodesik}{}{.}
Buktikan bahwa
\definitionsverweis {transpor paralel}{}{}
sepanjang $\gamma$ memetakan vektor tangensial

\mavergleichskette
{\vergleichskette
{ \gamma'(a)
}
{ \in }{ T_{\gamma(a)}Y
}
{  }{
}
{    }{
}
{   }{
}
}
{}{}{}
ke vektor tangensial

\mavergleichskette
{\vergleichskette
{ \gamma'(b)
}
{ \in }{ T_{\gamma(b)}Y
}
{  }{
}
{    }{
}
{   }{
}
}
{}{}{.}

}
{

Menurut
Lemma 6.8 (Differentialgeometrie (Osnabrück 2023))
medan kecepatan $\gamma'$ merupakan suatu
\definitionsverweis {medan vektor paralel}{}{}
sepanjang $\gamma$. Oleh karena itu,
\definitionsverweis {transpor paralel}{}{}
$\Psi_\gamma$ sepanjang $\gamma$ memetakan vektor

\mavergleichskette
{\vergleichskette
{ \gamma'(a)
}
{ \in }{  T_{\gamma(a)} Y
}
{  }{
}
{    }{
}
{   }{
}
}
{}{}{}
ke vektor

\mavergleichskette
{\vergleichskette
{ \gamma'(b)
}
{ \in }{  T_{\gamma(b)} Y
}
{  }{
}
{    }{
}
{   }{
}
}
{}{}{}
.

 }



\inputaufgabeklausurloesung
{0}

{

}
{  }



\inputaufgabeklausurloesung
{3}

{

Apakah pemetaan
\maabbeledisp {\varphi} {S^1} {S^1
} {(x,y)} {( \betrag {  x  } ,y)
} {,}
\definitionsverweis {terdiferensialkan}{}{?}

}
{

Kita meninjau komposisi pemetaan-pemetaan

\mathdisp {\R \stackrel{ \theta }{\longrightarrow} S^1 \stackrel{\varphi}{\longrightarrow} S^1 \subset \R^2 \stackrel{p_1}{\longrightarrow} \R} { , }

dengan

\mavergleichskettedisp
{\vergleichskette
{ \theta(t)
}
{ =} {(\cos  t , \sin  t )
}
{ } {
}
{ } {
}
{ } {
}
}
{}{}{}
dan $p_1$ merupakan proyeksi pertama. Parametrisasi trigonometrik, inklusi
\zusatzklammer {suatu submanifold tertutup} {} {,}
dan proyeksi tersebut semuanya diferensiabel. Jika $\varphi$ diferensiabel, maka komposisi keseluruhannya juga harus diferensiabel. Akan tetapi, komposisi itu adalah
\maabbeledisp {} {\R} {\R
} { t } {\betrag {  \cos  t  }
} {.}
Pemetaan ini tidak diferensiabel di
\mathl{t= \pi/2}{} karena gradien dari kiri di titik tersebut adalah $-1$, sedangkan gradien dari kanan adalah $1$. Jadi, pemetaan $\varphi$ tidak diferensiabel.

 }



\inputaufgabeklausurloesung
{7}

{

Misalkan $M$ suatu manifold diferensiabel kompak dengan bentuk volume positif kontinu $\omega$. Buktikan bahwa

\mavergleichskettedisp
{\vergleichskette
{ \int_M \omega
}
{ <} { \infty
}
{ } {
}
{ } {
}
{ } {
}
}
{}{}{}.

}
{

Untuk setiap titik

\mavergleichskette
{\vergleichskette
{ P
}
{ \in }{ M
}
{  }{
}
{    }{
}
{   }{
}
}
{}{}{}
terdapat suatu lingkungan peta terbuka

\mavergleichskette
{\vergleichskette
{ P
}
{ \in }{ U
}
{  }{
}
{    }{
}
{   }{
}
}
{}{}{}
dan suatu pemetaan peta
\maabbdisp {\alpha} { U } { V
} {}
dengan

\mavergleichskette
{\vergleichskette
{ V
}
{ \subseteq }{ \R^n
}
{  }{
}
{    }{
}
{   }{
}
}
{}{}{}
terbuka, sedemikian rupa sehingga

\mavergleichskette
{\vergleichskette
{ \alpha^{-1 *} \omega
}
{ = }{ fdx_1 \wedge \ldots \wedge dx_n
}
{  }{
}
{    }{
}
{   }{
}
}
{}{}{}
berlaku dengan
\mathl{f}{} kontinu dan positif. Kita juga memperoleh suatu lingkungan terbuka

\mavergleichskette
{\vergleichskette
{ P
}
{ \in }{ U'
}
{ \subseteq }{ U
}
{    }{
}
{   }{
}
}
{}{}{,}
yang homeomorf dengan suatu bola terbuka

\mavergleichskette
{\vergleichskette
{ U'
}
{ \cong }{ B'
}
{ \subseteq }{ V
}
{    }{
}
{   }{
}
}
{}{}{}
dan kita dapat menganggap bahwa penutupan bola tersebut seluruhnya terletak di $V$. Bola tertutup itu tertutup dan terbatas; karena itu, fungsi kontinu $f$ terbatas pada bola tersebut dan dengan demikian juga pada $U'$. Maka
\mathl{\int_{U'} \omega}{} berhingga, dengan $U'$ suatu lingkungan terbuka dari $P$.

Himpunan-himpunan terbuka

\mavergleichskette
{\vergleichskette
{ U'
}
{ = }{ U'(P)
}
{  }{
}
{    }{
}
{   }{
}
}
{}{}{}
ini menutupi $M$. Karena kekompakan, terdapat suatu penutup berhingga

\mavergleichskettedisp
{\vergleichskette
{ M
}
{ =} { \bigcup_{i = 1} ^n U_i
}
{ } {
}
{ } {
}
{ } {
}
}
{}{}{}
dengan

\mavergleichskettedisp
{\vergleichskette
{ \int_{U_i } \omega
}
{ <} { \infty
}
{ } {
}
{ } {
}
{ } {
}
}
{}{}{.}
Karena kepositifan, diperoleh

\mavergleichskettedisp
{\vergleichskette
{ \int_{M} \omega
}
{ \leq} { \sum_{i = 1}^n  \left(  \int_{U_i } \omega  \right)
}
{ <} { \infty
}
{ } {
}
{ } {
}
}
{}{}{.}

 }



\inputaufgabeklausurloesung
{0}

{

}
{  }



\inputaufgabeklausurloesung
{7 (1+2+4)}

{

\aufzaehlungdreiabc{ Berikan suatu sifat topologis yang menunjukkan bahwa
\definitionsverweis {interval satuan terbuka}{}{}
dan
\definitionsverweis {lingkaran}{}{}
$S^1$ tidak
\definitionsverweis {homeomorfik}{}{.}
}{Buktikan bahwa setiap
$1$-\definitionsverweis {bentuk diferensial}{}{}
kontinu pada
\mathl{]0,1[}{}
bersifat
\definitionsverweis {eksak}{}{.}
}{ Berikan suatu
$1$-bentuk konkret pada $S^1$ yang
\definitionsverweis {tertutup}{}{}
tetapi tidak eksak.
}

}
{

a) Menurut
Satz 13.14 (Differentialgeometrie (Osnabrück 2023))
$S^1$, sebagai himpunan bagian tertutup dan terbatas dari $\R^2$, bersifat
\definitionsverweis {kompak}{}{,}
sedangkan interval terbuka tersebut tidak kompak.

b) Berlaku

\mavergleichskettedisp
{\vergleichskette
{ \omega
}
{ =} { hdt
}
{ } {
}
{ } {
}
{ } {
}
}
{}{}{}
dengan suatu fungsi kontinu
\maabbdisp {h} {]0,1[} {\R
} {.}
Menurut
Hauptsatz der Infinitesimalrechnung,
$h$ mempunyai suatu antiturunan $H$. Dalam bahasa bentuk diferensial, hal ini berarti

\mavergleichskettedisp
{\vergleichskette
{H'
}
{ =} { hdt
}
{ =} { \omega
}
{ } {
}
{ } {
}
}
{}{}{,}
sehingga $\omega$ eksak.

c) Misalkan
\mathl{S^1 \subset \R^2}{} diberikan, dengan koordinat pada $\R^2$ dinotasikan oleh
\mathkor {} {u} {dan} {v} {}
. Kita meninjau bentuk diferensial

\mavergleichskettedisp
{\vergleichskette
{ \omega
}
{ =} { udv-vdu
}
{ } {
}
{ } {
}
{ } {
}
}
{}{}{}
pada $S^1$. Karena $S^1$ berdimensi satu, bentuk ini tertutup. Untuk lintasan
\maabbeledisp {\gamma} {[0, 2 \pi]} {S^1
} { t } { \begin{pmatrix}  \cos  t  \\ \sin  t     \end{pmatrix}
} {,}
berlaku

\mavergleichskettedisp
{\vergleichskette
{ \int_\gamma \omega
}
{ =} { \int_0^{2 \pi} \gamma^* \omega  dt
}
{ =} { \int_0^{2 \pi}  \cos  t  \cdot \cos  t  +  \sin  t   \sin  t   dt
}
{ =} { \int_0^{2 \pi} 1 dt
}
{ =} { 2 \pi
}
}
{}{}{.}
Menurut
Korollar 23.3 (Differentialgeometrie (Osnabrück 2023))
$\omega$ tidak mungkin eksak.

 }



\inputaufgabeklausurloesung
{4}

{

Misalkan

\mavergleichskette
{\vergleichskette
{ U
}
{ \subseteq }{ \R^n
}
{  }{
}
{    }{
}
{   }{
}
}
{}{}{}
\definitionsverweis {terbuka}{}{}
dan
\maabb {f} { U } { \R
} {}
suatu
\definitionsverweis {fungsi yang terdiferensialkan secara kontinu}{}{}
dua kali. Buktikan bahwa

\mavergleichskettedisp
{\vergleichskette
{ d(df)
}
{ =} { 0
}
{ } {
}
{ } {
}
{ } {
}
}
{}{}{,}
dengan $d$ menyatakan
\definitionsverweis {turunan eksterior}{}{.}

}
{

Untuk suatu bentuk-$1$

\mavergleichskette
{\vergleichskette
{ \omega
}
{ = }{ \sum_{j = 1}^n g_jdx_j
}
{  }{
}
{    }{
}
{   }{
}
}
{}{}{}
dengan menggunakan

\mavergleichskette
{\vergleichskette
{ dx_i \wedge dx_j
}
{ = }{ -dx_j \wedge dx_i
}
{  }{
}
{    }{
}
{   }{
}
}
{}{}{}

\mavergleichskettealign
{\vergleichskettealign
{d \omega
}
{ =} { \sum_{j = 1 }^n  d(g_j dx_j)
}
{ =} { \sum_{j = 1 }^n   \left(  \sum_{i = 1}^n { \frac{   \partial  g_j   }{  \partial   x_i   } } dx_i \wedge dx_j  \right)
}
{ =} { \sum _{1 \leq i <j\leq n} \left(  { \frac{   \partial  g_j   }{  \partial   x_i   } } - { \frac{   \partial  g_i   }{  \partial   x_j   } }  \right) dx_i \wedge dx_j
}
{ } {
}
}
{}
{}{.}
Untuk suatu fungsi $f$ yang dua kali terdiferensialkan secara kontinu, berlaku

\mavergleichskette
{\vergleichskette
{ df
}
{ = }{ \sum_{j = 1}^n g_j dx_j
}
{  }{
}
{    }{
}
{   }{
}
}
{}{}{}
dengan
\definitionsverweis {turunan-turunan parsial}{}{}

\mavergleichskette
{\vergleichskette
{ g_j
}
{ = }{ { \frac{   \partial   f   }{  \partial   x_j   } }
}
{  }{
}
{    }{
}
{   }{
}
}
{}{}{,}
dan karena itu

\mavergleichskette
{\vergleichskette
{ d(df)
}
{ = }{ 0
}
{  }{
}
{    }{
}
{   }{
}
}
{}{}{}
menurut
Satz von Schwarz.

 }



\inputaufgabeklausurloesung
{3}

{

Deskripsikan
\definitionsverweis {pemetaan}{}{}
\maabbeledisp {\psi} { \Complex \setminus \{1\} } { \Complex
} {w} { { \frac{   (w+1)  { \mathrm i}  }{ 1-w  } }
} {}
dalam koordinat riil.

}
{

Kita menulis

\mavergleichskette
{\vergleichskette
{ w
}
{ = }{ a+b { \mathrm i}
}
{  }{
}
{    }{
}
{   }{
}
}
{}{}{.}
Maka

\mavergleichskettealign
{\vergleichskettealign
{  \psi ( a+b { \mathrm i}   )
}
{ =} { { \frac{   \left(  a+b { \mathrm i} +1  \right) { \mathrm i}      }{  1-  a-b { \mathrm i}      } }
}
{ =} { { \frac{     a { \mathrm i} +{ \mathrm i}  -b  }{  1-  a-b { \mathrm i}      } }
}
{ =} { { \frac{     \left(  a { \mathrm i} +{ \mathrm i}  -b  \right) \left(  1-  a+b { \mathrm i}  \right)  }{  \left(  1-  a-b { \mathrm i}  \right) \left(  1-  a+b { \mathrm i}  \right)      } }
}
{ =} { { \frac{     \left(  a+1 \right) \left(  1-a  \right) { \mathrm i}  -b^2 { \mathrm i} -b \left(  1-  a \right) -b\left(  a+1  \right)  }{  \left(  1-  a \right)^2 +b^2  } }
}
}
{
\vergleichskettefortsetzungalign
{ =} { { \frac{     \left(  1-a ^2-b^2  \right) { \mathrm i}   -2 b  }{  \left(  1-  a \right)^2 +b^2  } }
}
{ } {
}
{ } {}
{ } {}
}
{}{.}
Dalam koordinat riil, dengan demikian,

\mavergleichskettedisp
{\vergleichskette
{ \psi \begin{pmatrix}  a  \\b   \end{pmatrix}
}
{ =} { { \frac{   1  }{  -2a+1+a^2+b^2 } } \begin{pmatrix}  -2b \\ 1-a^2-b^2   \end{pmatrix}
}
{ } {
}
{ } {
}
{ } {
}
}
{}{}{.}

 }



\inputaufgabeklausurloesung
{3}

{

Misalkan

\mavergleichskette
{\vergleichskette
{ H
}
{ \subseteq }{ \R^n
}
{  }{
}
{    }{
}
{   }{
}
}
{}{}{}
suatu
\definitionsverweis {setengah ruang Euklides}{}{}
dan

\mavergleichskette
{\vergleichskette
{ P
}
{ \in }{ H
}
{  }{
}
{    }{
}
{   }{
}
}
{}{}{.}
Andaikan terdapat suatu himpunan yang terbuka di $\R^n$, yaitu $V$, dengan

\mavergleichskette
{\vergleichskette
{ P
}
{ \in }{ V
}
{ \subseteq }{ H
}
{    }{
}
{   }{
}
}
{}{}{.}
Buktikan bahwa
\mathl{P}{} bukan titik batas dari $H$.

}
{

Misalkan

\mavergleichskettedisp
{\vergleichskette
{ P
}
{ =} { \begin{pmatrix}  x_1  \\ x_2 \\  \vdots\\ x_n   \end{pmatrix}
}
{ } {
}
{ } {
}
{ } {
}
}
{}{}{.}
Kita dapat mengganti, di $\R^n$, lingkungan terbuka $V$ dengan suatu bola terbuka
\zusatzklammer {di $\R^n$} {} {}
\mathl{U \left(   P , \epsilon  \right)}{} dengan

\mavergleichskette
{\vergleichskette
{ P
}
{ \in }{   U \left(   P , \epsilon  \right)
}
{ \subseteq }{  V
}
{ \subseteq   }{ H
}
{   }{
}
}
{}{}{}
. Jika $P$ suatu titik batas, maka

\mavergleichskette
{\vergleichskette
{ x_1
}
{ = }{ 0
}
{  }{
}
{    }{
}
{   }{
}
}
{}{}{.}
Akan tetapi, dalam hal itu

\mavergleichskette
{\vergleichskette
{ \begin{pmatrix}  - { \frac{   \epsilon }{  2 } }  \\ x_2 \\  \vdots\\ x_n   \end{pmatrix}
}
{ \in }{ U \left(   P , \epsilon  \right)
}
{  }{
}
{    }{
}
{   }{
}
}
{}{}{,}
padahal titik ini tidak termasuk dalam $H$.

 }



\inputaufgabeklausurloesung
{7}

{

Misalkan $M$ suatu
\definitionsverweis {manifold diferensiabel}{}{}
berdimensi $n$ dan

\mavergleichskette
{\vergleichskette
{ P
}
{ \in }{ M
}
{  }{
}
{    }{
}
{   }{
}
}
{}{}{}
suatu titik. Buktikan bahwa terdapat suatu
\definitionsverweis {manifold dengan batas}{}{}
$N$ yang batasnya $\partial N$ difeomorfik dengan permukaan bola berdimensi $(n-1)$, yaitu $S^{n-1}$, sedemikian sehingga
\mathkor {} {M \setminus \{ P \}} {dan} {N \setminus \partial N} {}
saling difeomorfik.

}
{

Misalkan

\mavergleichskette
{\vergleichskette
{ P
}
{ \in }{  U
}
{ \subseteq }{  M
}
{    }{
}
{   }{
}
}
{}{}{}
suatu daerah peta terbuka dengan peta
\maabbdisp {\alpha} { U } {V \subseteq \R^n
} {.}
Kita dapat menganggap bahwa $\alpha(P)$ adalah titik nol dan $V$ adalah bola terbuka berjari-jari $3$. Misalkan
\maabbdisp {\varphi} { U \left(   0 , 3  \right) \setminus \{0\}  } { U \left(   0 , 3  \right) \setminus B  \left(  0 , 1 \right)
} {}
suatu difeomorfisme yang merupakan identitas pada
\mathl{U \left(   0 , 3  \right) \setminus U \left(   0 , 2  \right)}{} dan memetakan
\mathl{U \left(   0 , 2  \right) \setminus \{0\}}{} ke
\mathl{U \left(   0 , 2  \right) \setminus B  \left(  0 , 1 \right)}{}. Pemetaan seperti ini diperoleh dengan meregangkan titik-titik menggunakan fungsi
\maabbdisp {h} {[0,3]} { [0,3]
} {}
dengan

\mavergleichskettedisp
{\vergleichskette
{ h(r)
}
{ =} { \begin{cases} 1+ { \frac{   1  }{  4 } } r^2 \text{ jika } 0 \leq r \leq 2 \, , \\ r \text{ jika } r >  2  \, , \end{cases}
}
{ } {
}
{ } {
}
{ } {
}
}
{}{}{}
. Selanjutnya, citra difeomorfisme $\varphi$, yaitu
\mathl{U \left(   0 , 3  \right) \setminus B  \left(  0 , 1 \right)}{,} dapat dipandang sebagai

\mavergleichskettedisp
{\vergleichskette
{  U \left(   0 , 3  \right) \setminus B  \left(  0 , 1 \right)
}
{ \subset} {  U \left(   0 , 3  \right) \setminus U \left(   0 , 1  \right)
}
{ } {
}
{ } {
}
{ } {
}
}
{}{}{}
. Himpunan yang lebih besar merupakan suatu manifold dengan batas $S \left(   0 ,  1  \right)$. Misalkan $N$ manifold yang diperoleh dengan mengganti $U$ oleh
\mathl{U \left(   0 , 3  \right) \setminus U \left(   0 , 1  \right)}{}. Permukaan bola tersebut menjadi batas dari $N$. Difeomorfisme $\varphi$ dapat diperluas menjadi suatu difeomorfisme
\maabbdisp {} {M \setminus \{P\} } { N \setminus \partial N
} {}
karena pada daerah transisi terbuka
\mathl{U \left(   0 , 3  \right) \setminus U \left(   0 , 2  \right)}{} pemetaan tersebut merupakan identitas.

 }



\inputaufgabeklausurloesung
{2}

{

Berikan suatu
\definitionsverweis {penghabisan kompak}{}{}
untuk $\R^d$.

}
{

Kita meninjau
\definitionsverweis {bola-bola tertutup}{}{}

\mathbed {B  \left(  0 , n  \right)} {}
 {n \in \N_+} {}
 {} {} {} {,}
yang
\definitionsverweis {kompak}{}{}
dan menutupi $\R^d$. Interior terbuka dari
\mathl{B  \left(  0 , n+1  \right)}{}
adalah
\definitionsverweis {bola terbuka}{}{}

\mathl{U \left(   0,n+1  \right)}{}, dan karena

\mavergleichskettedisp
{\vergleichskette
{ B  \left(  0,n \right)
}
{ \subseteq} { U \left(   0,n+1  \right)
}
{ } {
}
{ } {
}
{ } {
}
}
{}{}{}
diperoleh suatu
\definitionsverweis {penghabisan kompak}{}{.}

 }



\inputaufgabeklausurloesung
{9}

{

Buktikan teorema titik tetap Brouwer.

}
{

Untuk menyederhanakan notasi, misalkan

\mavergleichskette
{\vergleichskette
{ r
}
{ = }{ 1
}
{  }{
}
{    }{
}
{   }{
}
}
{}{}{.}  Andaikan terdapat suatu pemetaan $\psi$ tanpa titik tetap yang terdiferensialkan secara kontinu. Maka selalu

\mavergleichskettedisp
{\vergleichskette
{ x
}
{ \neq} { \psi(x)
}
{ } {
}
{ } {
}
{ } {
}
}
{}{}{,}
sehingga kedua titik tersebut menentukan suatu garis. Gagasannya adalah menggunakan garis ini untuk mengambil salah satu
\zusatzklammer {dari kedua} {} {}
titik potong dengan permukaan bola sebagai titik citra suatu retraksi ke batas. Dengan fungsi bantu

\mavergleichskettedisp
{\vergleichskette
{ h(x)
}
{ =} { { \frac{   \psi(x)-x }{  \Vert { \psi (x)-x} \Vert  } }
}
{ } {
}
{ } {
}
{ } {
}
}
{}{}{}
kita mendefinisikan suatu pemetaan
\maabbdisp {\varphi} { B  \left(  0 , 1 \right) } { \R^n
} {}
melalui

\mavergleichskettedisphandlinks
{\vergleichskettedisphandlinks
{ \varphi(x)
}
{ =} { x + \left(  - \left\langle  x  , h(x)  \right\rangle + \sqrt{1 + \left\langle  x  , h(x)  \right\rangle^2 - \Vert { x} \Vert^2  }  \right) \cdot h(x)
}
{ } {
}
{ } {
}
{ } {
}
}
{}{}{.}
Ekspresi di bawah akar tersebut positif. Hal ini jelas untuk

\mavergleichskette
{\vergleichskette
{ \Vert { x } \Vert
}
{ < }{ 1
}
{  }{
}
{    }{
}
{   }{
}
}
{}{}{}
dan untuk

\mavergleichskette
{\vergleichskette
{ \Vert { x } \Vert
}
{ = }{ 1
}
{  }{
}
{    }{
}
{   }{
}
}
{}{}{}
diperoleh suatu titik pada permukaan bola yang garis penghubungnya dengan titik bola
\mathl{\psi(x)}{} tidak tegak lurus terhadap $x$
\zusatzklammer {ruang singgung afin pada suatu titik permukaan bola hanya berpotongan dengan bola di satu titik} {} {,}
sehingga

\mavergleichskettedisp
{\vergleichskette
{ \left\langle  x  , h(x)  \right\rangle
}
{ \neq} { 0
}
{ } {
}
{ } {
}
{ } {
}
}
{}{}{}
berlaku. Karena akar kuadrat dan nilai mutlak di luar titik nol
\definitionsverweis {terdiferensialkan secara kontinu}{}{,}
baik
\mathkor {} {h(x)} {maupun} {\varphi(x)} {}
merupakan pemetaan yang terdiferensialkan secara kontinu. Menurut
Aufgabe 23.23 (Differentialgeometrie (Osnabrück 2023))
pemetaan $\varphi$ memetakan bola ke permukaan bola dan pembatasannya pada permukaan bola adalah identitas. Dengan demikian, diperoleh suatu
\definitionsverweis {retraksi}{}{}
yang terdiferensialkan secara kontinu dari bola penuh tertutup ke batasnya, yang menurut
Satz 23.9 (Differentialgeometrie (Osnabrück 2023))
tidak mungkin ada.

 }



\inputaufgabeklausurloesung
{7 (1+2+4)}

{

Pada
\definitionsverweis {bundel vektor trivial}{}{}

\mathl{\R \times \R^2}{} di atas $\R$, kita meninjau
\definitionsverweis {koneksi trivial}{}{}
$\nabla$.
\aufzaehlungdrei{Buktikan bahwa

\mavergleichskette
{\vergleichskette
{ f_1(t)
}
{ = }{ \left(  1+ t^2  , \,  t^3                   \right)
}
{  }{
}
{    }{
}
{   }{
}
}
{}{}{}
dan

\mavergleichskette
{\vergleichskette
{ f_2(t)
}
{ = }{ \left(  t  , \,   1+t^2                   \right)
}
{  }{
}
{    }{
}
{   }{
}
}
{}{}{}
merupakan
\definitionsverweis {penampang basis}{}{}
dari bundel vektor tersebut.
}{Nyatakan penampang basis standar $e_1,e_2$ sebagai kombinasi linear dari penampang basis $f_1,f_2$.
}{Tentukan
\definitionsverweis {simbol-simbol Christoffel}{}{}
dari $\nabla$ terhadap penampang-penampang basis tersebut.
}

}
{

\aufzaehlungdrei{Berlaku

\mavergleichskettedisp
{\vergleichskette
{ \det  \begin{pmatrix}  1+t^2   &  t^3  \\  t  &  1+t^2  \end{pmatrix}
}
{ =} { 1+2t^2+t^4 -t^4
}
{ =} { 1 +2t^2
}
{ >} {  0
}
{ } {
}
}
{}{}{}
selalu positif. Oleh karena itu, kedua vektor
\mathkor {} {f_1(t)} {dan} {f_2(t)} {}
untuk setiap

\mavergleichskette
{\vergleichskette
{ t
}
{ \in }{  \R
}
{  }{
}
{    }{
}
{   }{
}
}
{}{}{}
membentuk suatu basis.
}{Dari hubungan

\mavergleichskettedisp
{\vergleichskette
{ \begin{pmatrix} f_1  \\f_2   \end{pmatrix}
}
{ =} { \begin{pmatrix}  1+t^2  &   t^3  \\  t  &  1+t^2 \end{pmatrix}  \begin{pmatrix} e_1  \\e_2   \end{pmatrix}
}
{ } {
}
{ } {
}
{ } {
}
}
{}{}{}
diperoleh

\mavergleichskettedisp
{\vergleichskette
{ \begin{pmatrix} e_1  \\e_2   \end{pmatrix}
}
{ =} { { \frac{   1  }{  1+2t^2 } } \begin{pmatrix}  1+t^2  &   -t^3  \\  -t &  1+t^2 \end{pmatrix}  \begin{pmatrix} f_1  \\f_2   \end{pmatrix}
}
{ } {
}
{ } {
}
{ } {
}
}
{}{}{,}
jadi

\mavergleichskettedisp
{\vergleichskette
{ e_1
}
{ =} { { \frac{   1+t^2 }{  1+2t^2 } }   f_1 - { \frac{  t^3 }{  1+2t^2 } } f_2
}
{ } {
}
{ } {
}
{ } {
}
}
{}{}{}
dan

\mavergleichskettedisp
{\vergleichskette
{ e_2
}
{ =} { -{ \frac{   t  }{  1+2t^2 } }   f_1 + { \frac{   1+t^2 }{  1+2t^2 } } f_2
}
{ } {
}
{ } {
}
{ } {
}
}
{}{}{.}
}{Untuk koneksi trivial, berlaku

\mavergleichskettedisp
{\vergleichskette
{ \nabla_\partial (f)
}
{ =} { \partial f
}
{ } {
}
{ } {
}
{ } {
}
}
{}{}{.}
Dengan demikian,

\mavergleichskettealign
{\vergleichskettealign
{ \nabla_\partial (f_1)
}
{ =} { \nabla_\partial \left(  \left(  1+t^2  \right) e_1 + t^3e_2  \right)
}
{ =} { 2t e_1 + 3t^2 e_2
}
{ =} { 2t  \left(  { \frac{   1+t^2 }{  1+2t^2 } }   f_1 - { \frac{  t^3 }{  1+2t^2 } } f_2   \right)  + 3t^2  \left(  -{ \frac{   t  }{  1+2t^2 } }   f_1 + { \frac{   1+t^2 }{  1+2t^2 } } f_2   \right)
}
{ =} {  \left(   { \frac{   2t+2t^3 }{  1+2t^2 } } - { \frac{   3t^3 }{  1+2t^2 } }  \right) f_1  +  \left(  - { \frac{  t^3 }{  1+2t^2 } }   + { \frac{   3t^2+3t^4 }{  1+2t^2 } }     \right)   f_2
}
}
{
\vergleichskettefortsetzungalign
{ =} { { \frac{   2t - t^3 }{  1+2t^2 } }  f_1  +  { \frac{   3t^2-t^3+3t^4 }{  1+2t^2 } }      f_2
}
{ } {}
{ } {}
{ } {}
}
{}{}
dan

\mavergleichskettealign
{\vergleichskettealign
{ \nabla_\partial (f_2)
}
{ =} { \nabla_\partial \left(  t e_1 + \left(  1+t^2  \right) e_2  \right)
}
{ =} { e_1 + 2t  e_2
}
{ =} { { \frac{   1+t^2 }{  1+2t^2 } } f_1 - { \frac{  t^3 }{  1+2t^2 } } f_2 + 2t \left(  -{ \frac{   t  }{  1+2t^2 } } f_1 + { \frac{   1+t^2 }{  1+2t^2 } } f_2   \right)
}
{ =} { \left( { \frac{   1+t^2 }{  1+2t^2 } }  - { \frac{   2t^2 }{  1+2t^2 } }  \right) f_1  + \left(  - { \frac{   t^3  }{  1+2t^2  } } +  { \frac{   2t+2t^3  }{  1+2t^2 } }  \right) f_2
}
}
{
\vergleichskettefortsetzungalign
{ =} { { \frac{   1 - t^2  }{  1+2t^2 } } f_1 + { \frac{   2t   +t^3 }{  1+2t^2 } }      f_2
}
{ } {}
{ } {}
{ } {}
}
{}{.}
Fungsi-fungsi koefisien tersebut adalah simbol-simbol Christoffel.
}

 }
